% !TEX root = ../main.tex
\section{问题分析}

\subsection{问题一的分析}

一次带有 $\pm1^\circ$ 误差的示向观测对应由两条边界射线确定的角形区域，该角形可写成两个有向半平面的交。将全部观测对应的半平面取交，可得到凸多边形定位区域。由此，问题一可分为三个连续步骤：先利用半平面交算法计算多边形顶点，再利用旋转卡壳算法求最远顶点对，最后检查以最远顶点对为直径的圆是否包含全部多边形顶点。

\subsection{问题二的分析}

第二检测点的选择同时受到接收距离和交会角的限制。若第二点仍位于第一次示向线附近，则两次观测近乎平行，角度误差会被显著放大；若只追求接近 $90^\circ$ 的交会角，又可能使第二点超出干扰源的有效接收范围。为此，首先利用目标圆域、第一次示向角和 $1500$ 米接收上界构造候选源区域，再以 $1000$ 米接收下界构造具有确定性接收保证的检测点区域。在该区域的左右两个子区域中，以连续最坏定位损失为理论主目标，以移动时间为次目标；数值求解则明确报告有限样本和有限候选点网格下的近似结果。

\subsection{问题之间的联系}

问题一给出了由多次示向观测更新定位区域并计算其直径的方法；问题二将定位区域直径作为检测点质量的直接评价指标，并反向决定下一次检测位置。因此，问题一是几何计算基础，问题二是建立在该基础上的主动测向决策模型。
